\section{The optimal approximation and its scalar consequences}
\label{p13:optimal:section}

The arithmetic and probabilistic estimates now meet at a
single quantity, the absolute third-order activity.
We complete the main theorem, then record what its precision
allows for unbounded scalar observables and for the choice
of mean.

\begin{proof}[Proof of Theorem~\ref{p1:thm:main}]
Write $\varepsilon_{\mathrm{out}}=\varepsilon$ for the requested
final margin and choose
$0<\tau<\min(\varepsilon_{\mathrm{out}}/4,1/960)$.
Invoke Theorem~\ref{p13:transfer:reconstruction} with accuracy $\tau$
and Corollary~\ref{p13:third:activity} with accuracy $\tau$.
The triple contribution is
\[
 N^\tau A_3\ll N^{-41/30+2\tau}
       =N^{-4/3-(1/30-2\tau)}.
\]
Since $1/30-2\tau>0$ and the other reconstruction term is
$N^{-4/3+\tau}$, this proves the upper bound.
Theorem~\ref{p13:lower:optimal} gives the lower bound for the same
exact parameters, using its own fixed Charlier interpolation order.

The mean estimate is Proposition~\ref{p1:arith:core}.
The pair asymptotic \eqref{p1:eq:alpha-leading} shows
that \(\alpha_N>0\), \(\alpha_N\to0\), and \(\mu_N\)
stays in a compact subset of \((0,\infty)\).
Lemma~\ref{p1:close:positivity} therefore makes the
reference a probability measure for sufficiently large \(N\).
Every auxiliary accuracy parameter and fixed order was chosen after
\(C,\varepsilon\) and before \(N\). The conclusion
is uniform on the entire fixed critical band.
\end{proof}

\subsection{Polynomial weights and fixed moments}

\begin{corollary}[Every fixed polynomial weight]
\label{p13:optimal:weighted}
For every fixed integer \(m\ge0\) and \(\varepsilon>0\),
\begin{equation}
 \sum_{k\ge0}(1+k)^m
   \left|\P(Z_N=k)-Q_{\mu_N,\alpha_N}\{k\}\right|
       \ll_{C,m,\varepsilon}N^{-4/3+\varepsilon}.
 \label{p13:optimal:weighted-bound}
\end{equation}
For every fixed \(r\ge2\), consequently,
\begin{equation}
 \E(Z_N)_r
 =\mu_N^r+\binom r2\mu_N^{r-2}\alpha_N
           +O_{C,r,\varepsilon}(N^{-4/3+\varepsilon}).
 \label{p13:optimal:moments}
\end{equation}
The equalities for \(r=0,1,2\) hold without a remainder
when the right side is interpreted in its natural form.
\end{corollary}

\begin{proof}
Let \(R=\Law Z_N-Q_{\mu_N,\alpha_N}\).
For every fixed integer \(v>1\), H\"older gives
\[
 \sum_k(1+k)^m|R\{k\}|
 \le\|R\|_1^{1-1/v}
       \left(\sum_k(1+k)^{mv}|R\{k\}|\right)^{1/v}.
\]
The second factor is uniformly bounded.
For the true law this follows from the independent
fixed-moment proof in Theorem~\ref{p1:fixed:moments};
for the reference it follows by dominating its variation
by a fixed sum of shifted Poisson laws, as in
Theorem~\ref{p13:lower:optimal}.
Choose \(v\) sufficiently large after \(\varepsilon\)
and then a sufficiently small margin in the main theorem.
This proves \eqref{p13:optimal:weighted-bound}.
The factorial generating function of \(Q_{\mu,\alpha}\)
is \(e^{\mu t}(1+\alpha t^2/2)\); its coefficient of
\(t^r\), multiplied by \(r!\), gives the moment formula.
\end{proof}

This interpolation concerns the scalar count and fixed
polynomial degrees. The positional comparison in
Lemma~\ref{p1:limits:vector-reference} has its own
vector expansion and retains its
\(N^{-4/5+\varepsilon}\) error.

\begin{corollary}[No retained start]
\label{p13:optimal:zero}
Uniformly in a fixed critical band,
\[
 \P(Z_N=0)
   =e^{-\mu_N}\left(1+\frac{\alpha_N}{2}\right)
          +O_{C,\varepsilon}(N^{-4/3+\varepsilon}).
\]
\end{corollary}

\begin{proof}
Evaluate the approximation at zero, using
\((\Pi_\mu*\Delta^{*2})\{0\}=e^{-\mu}\).
\end{proof}

The event is the absence of a retained start in \(I_N\).
It does not exclude a long run that started before \(N\).

\subsection{Why the mean must be specified at this precision}

\begin{proposition}[The quantized mean bias]
\label{p13:optimal:nominal}
Let \(\eta_x\) indicate one-window compatibility and put
\(h_x=L-\dim S_x\). Then
\begin{equation}
 \mu_N-\lambda_N=pM_{N,L},\qquad
 M_{N,L}=\sum_{x\in I_N}(\eta_x2^{h_x}-1)\in\Z.
 \label{p13:optimal:integer-mean}
\end{equation}
Thus either \(\mu_N=\lambda_N\) or
\(|\mu_N-\lambda_N|\ge p\asymp_C N^{-1}\).
Keeping the same exact \(\alpha_N\), one has
\[
 \|\Law Z_N-Q_{\lambda_N,\alpha_N}\|_1
       \ll_{C,\varepsilon}N^{-1+\varepsilon}.
\]
On every sufficiently large admissible pair \((N,L)\)
for which \(M_{N,L}\ne0\),
\begin{equation}
 \|\Law Z_N-Q_{\lambda_N,\alpha_N}\|_1
       \asymp_C|\mu_N-\lambda_N|.
 \label{p13:optimal:nominal-distance}
\end{equation}
No assertion that \(M_{N,L}\ne0\) at infinitely many
scales is included.
\end{proposition}

\begin{proof}
Sum the one-window formula
\(\E J_x=\eta_xp2^{h_x}\) over the exactly \(N\)
centres to obtain \eqref{p13:optimal:integer-mean}.
For two means in the compact range in question,
\eqref{p1:close:reference-lipschitz} gives
\[
 \|Q_{\mu,\alpha_N}-Q_{\lambda,\alpha_N}\|_1
 \ll_C|\mu-\lambda|.
\]
The reverse bound follows from the mass at zero,
\(e^{-\theta}(1+\alpha_N/2)\), by the mean value theorem.
The upper bound for the nominal reference follows
from the main theorem and
\(|\mu_N-\lambda_N|\ll N^{-1+\varepsilon}\).

If \(M_{N,L}\ne0\), the mean difference is at least \(p\).
Choose a fixed margin less than \(1/3\) in the main
theorem; its remainder is then \(o(p)\).
The triangle and reverse triangle inequalities transfer
the two-sided reference comparison to
\eqref{p13:optimal:nominal-distance}.
\end{proof}

Irregular compatible and incompatible windows contribute
with opposite signs to \(M_{N,L}\). Their presence alone
does not prove a nonzero total bias.
